% Generated draft: requires mathematical and source review.
\documentclass[11pt]{article}
\usepackage[margin=25mm]{geometry}
\usepackage{fontspec}
\setmainfont{Latin Modern Roman}
\usepackage{amsmath,amssymb,mathtools,unicode-math}
\setmathfont{Latin Modern Math}
\usepackage{enumitem,needspace,xurl,hyperref}
\setlength{\parindent}{0pt}
\setlength{\parskip}{5pt}
\setlength{\emergencystretch}{3em}
\newcommand{\N}{\mathbb{N}}
\newcommand{\Z}{\mathbb{Z}}
\newcommand{\Q}{\mathbb{Q}}
\newcommand{\R}{\mathbb{R}}
\newcommand{\C}{\mathbb{C}}
\newcommand{\F}{\mathbb{F}}
\newcommand{\A}{\mathbb{A}}
\newcommand{\PP}{\mathbb{P}}
\newcommand{\E}{\mathbb{E}}
\newcommand{\eps}{\varepsilon}
\newcommand{\dd}{\,\mathrm{d}}
\begin{document}
\section*{UnsolvedMath ID: 5100042}
\textbf{Source problem code:} catalogue entry
\subsection{Definitions and mathematical statement}
Fix a nondegenerate ellipse \(E\subset\mathbb R^2\), with center \(O\) and foci \(f_1,f_2\), and a confocal caustic \(C\). Let \(\mathcal F\) be a Poncelet family of nonsingular \(N\)-periodic billiard polygons \(P=(P_1,\ldots,P_N)\) on \(E\), whose edges are tangent to \(C\), in their billiard traversal order. The assertion is for every \(N\) admitting such a family; it does not assert that a fixed caustic supports every period. Put \(P_{N+1}=P_1\).

For \(j\in\{1,2\}\), unit-circle inversion about the focus \(f_j\) is
\[
 I_j(x)=f_j+\frac{x-f_j}{|x-f_j|^2}\qquad(x\ne f_j).
\]
The foci of a nondegenerate ellipse lie off its boundary, so the inverted vertices \(P_{j,i}^{\dagger}=I_j(P_i)\) are defined. Join them in the original traversal order to form \(P_j^{\dagger}\), and define its perimeter, including for a self-intersecting polygon, by
\[
 L_j^{\dagger}(P)=\sum_{i=1}^N
 \left|I_j(P_{i+1})-I_j(P_i)\right|.
\]
Invariant \(k_{802}\) states
\[
 \forall j\in\{1,2\}\ \exists c_j(E,C,N,\mathcal F)\
 \forall P\in\mathcal F:\qquad L_j^{\dagger}(P)=c_j(E,C,N,\mathcal F).
\]
Thus the ellipse, caustic, period, family, and inversion radius are fixed while the orbit varies. No equality of the two focal constants or formula for them is requested.

The source also introduces outer tangent-intersection and inner tangency-point polygons, various areas, angles, and polar/dual constructions for its invariant table; they do not occur in \(k_{802}\) and impose no additional condition on the perimeter assertion. The precise admissible period and caustic-degeneracy conventions are not stated in the excerpt, so the formulation is restricted to families where the described billiard polygons are defined.

The registry labels this solved and attributes the all-\(N\) perimeter result to Roitman--Garcia--Reznik, \emph{AMJ} (2021), Corollary 1(iii), also mentioning an \(N=3\) proof in arXiv:2012.03020. These are supplied literature attributions; the original invariant list's open label has no recorded later authoritative resolution check.
\subsection{Short English statement}
For any fixed Poncelet family of periodic billiard orbits in an ellipse, invert every orbit vertex in the unit circle about either focus. The perimeter of the resulting ordered polygon should remain constant as the orbit varies. The supplied summary reports this invariant proved for all admissible periods.
\subsection{Sources}
\begin{itemize}
\item Local registry: \url{data/unsolvedmath/problems.json}, permanent UnsolvedMath ID \(5100042\), source code catalogue entry, invariant \(k_{802}\); supplied statement, background, and research summary, updated 2026-08-06.
\item Dan Reznik, Ronaldo Garcia, and Jair Koiller, \emph{Eighty New Invariants of N-Periodics in the Elliptic Billiard} (2021), table row \(k_{802}\): \url{https://arxiv.org/abs/2004.12497}. Supplied source-TeX extraction accessed 2026-07-29.
\item Roitman--Garcia--Reznik, \emph{AMJ} (2021), Corollary 1(iii), as cited by the supplied research summary; no full title or publication details are supplied.
\item arXiv:2012.03020, cited by the summary for the \(N=3\) case; no title or authors are included in that reference in the record.
\end{itemize}
\end{document}
